% Time form: Present & past tense

\chapter{Discussion}\label{ch:discussion}

This chapter connects and interprets the findings of the Chapter~\ref{ch:evaluation},
with respect to the evaluation goals defined in
Section~\ref{subsec:evaluation-goals}.
Especially compared is the effectiveness of the \gls{AI}-based reconnaissance scenario
(Section~\ref{subsubsec:ai-scenario-evaluation}) and its
comparison against the manually engineered reference scenario
(Section~\ref{subsec:ai-vs-engineered-scenarios}).
While working with \gls{AGENTIC-AI} the results can not be separated from a model evaluation.
The summary draws the overall picture across both evaluation environments.
The subsequent sections then discuss the virtual \gls{Containerlab} and the real Cyberlab BFH
segment separately, because the two environments differ in scale and complexity.
On the one hand, the BFH Network Security Lab must not be exposed,
but it is also a black box when it comes to the correctness of its findings.
Therefore we evaluate the successful runs of the scenarios and compare the effectiveness of the agent to gather information.

\section{Summary}\label{sec:discussion-summary}

The evaluation shows that \gls{AI}-assisted reconnaissance is feasible but strongly
dependent on the underlying model.
The frontier model \textit{claude-opus-4-7} delivers the most convincing results
across both environments.
In the rather small virtual \gls{Containerlab} environment it reaches the highest correctness (9)
and lowest hallucination (9) scores.
Next to all discovered findings the detailed report on how to deal with the flaws is really helpful for further
investigation and a first lead on how to treat the vulnerabilities.
Also helpful is the extraction of some artifacts like cleartext passwords and credentials and further flaws in systems.
The two smaller models are delivering worse results.
The \textit{gpt-oss:120b} is adequate but prone to fabricating \gls{CVE} identifiers,
\textit{qwen3:30b} performs host and service
discovery correctly but lacks the enumeration depth for Critical finding.
The implemented static scenario provides the highest success rate, with 5 hosts, all 13 services and a fast execution speed of 29\,s.
No reasoning and evaluation is provided and leaves the operator with the task to assess the results.
But the rapid speed and not existing token costs make up for it.
It retrieves the \gls{LDAP} naming context yet never performs the anonymous bind that would expose the credentials.
This captures the central trade-off of the comparison: the engineered scenario offers
holistic and deterministic coverage, whereas the \gls{AI} scenario offers more in-depth analysis.

Multi-agent task decomposition improves the usability of smaller models without harming the frontier model.
\textit{qwen3:30b} cuts its token usage by roughly 68\,\% (86,446 to 27,789) and its
duration by ca 30\,\%.
\textit{Claude-opus-4-7} keeps identical discovery quality with slightly more overhead (+13\,\% tokens, 122\,s to 150\,s).
Decomposition also raises execution reliability.
The success factor improves for all three models, with \textit{claude-opus-4-7} reaching a perfect 10/10.

The picture transfers to the real Cyberlab BFH segment, a network roughly seven times larger than
the \gls{Containerlab} with about 35 live hosts on the 193.5.80.0/24 subnet.
There the multi-agent \textit{claude-opus-4-7} setup discovers on average 28 hosts, 53 services, and
52 findings per run and consistently surfaces the six high-severity structural risks.
The unstructured variant of the same model reaches an even higher correctness (9 vs.\ 7) because it
issues additional \gls{UDP} scans and \gls{NSE} scripts, but only at a 5.5\,\% success rate
and roughly three times the token cost (222,094 vs.\ 81,453).
The locally hosted \textit{qwen3:8b} fails completely with zero tool calls, and the static scenario
needs 1780\,s to enumerate the whole network.
The real environment therefore confirms the \gls{Containerlab} picture: the frontier model scales,
the engineered scenario stays reproducible but slow, and reliability remains the main open issue.


\section{Virtual Environment (Containerlab)}\label{sec:discussion-containerlab}

\paragraph{AI versus engineered reconnaissance.}

The comparison against the static scenario answers the goal (Section~\ref{subsec:ai-vs-engineered-scenarios}).
On just the level of reproducibility and raw coverage the engineered scenario is unbeatable.
It executes a fixed drill order and returns an identical structured \textit{NetworkDiscoveryResultMap}
on every run.
All services are discovered issued by an unconditional port sweep, where the \gls{AI} models compress the services to 7-9
or terminate scanning in a halfway finished process.

\todo{Add this part of the thesis}

However, raw coverage is not an indicator of intelligence.
The generalistic way how the engineered scenario was implemented, prevent the extraction of the cleartext credentials,
does not elaborate the printer's banner mismatch, and produces no severity-rated assessment.
The time needed to implement the right attack vector to access the vulnerability of different systems or services, is immense.
The frontier \gls{AI} model, despite enumerating fewer services, delivers the most holistic report and findings
by chaining observations into an attack context (credential reuse → lateral movement via SSH or SMB).
This addresses the creativity criteria from Section~\ref{subsec:criteria}: the \gls{AI} scenario identifies connections
that the engineered scenario, by construction, cannot.

\paragraph{Single-agent versus multi-agent.}

The two configurations confirm that smaller models benefit more from decomposition, where the larger models are not
influenced or less influenced by bigger context.
A frontier model already manages a large context well, so decomposition mainly adds orchestration overhead in the \gls{Containerlab}.

Smaller models are more failure prone as their context fills with accumulated tasks and history.
The smaller tasks and focused context restore both efficiency and token reduction for the local hosted model and a
more consistency that reflects within a tighter scatter cluster in Figure~\ref{fig:benchmark_scatter_containerlab}.
This is the primary architectural and conceptual finding which motivates to use the multi-agent setup as the default configuration.


\section{Real Network (Cyberlab BFH)}\label{sec:discussion-cyberlab}

\paragraph{Scaling to a real network.}

The Cyberlab BFH segment validates the approach on a real network roughly seven
times larger than \gls{Containerlab} (\(\approx\)35 live hosts vs.\ 5).
The multi-agent setup scaled to this size, discovering on average 28 hosts,
53 services, and 52 findings per run, and consistently surfaced the six
high-severity structural risks.
Unlike the \gls{Containerlab}, the lab is a black box without a published ground truth and must
not be exposed, so the discussion is limited to the successful runs and the effectiveness of the
agent to gather information rather than a full per-finding comparison.
In this more complex environment the frontier model also benefits from decomposition through a
higher success rate, where 5 of 27 multi-agent attempts produced a fully structured result.

\paragraph{Structured multi-agent versus unstructured reasoning.}

The real network exposes a trade-off that the virtual environment does not.
The unstructured-output configuration reaches a higher correctness (9 vs.\ 7) than the structured
multi-agent setup, because it issues \gls{UDP} scans alongside \gls{TCP} and executes
\gls{NSE} scripts that the specialised sub-agents omit, surfacing all three discriminating
findings (\gls{SNMP}, open \gls{DNS} recursion, and the cloned \gls{SSH} host keys).
The schema constraints that help the multi-agent setup with speed and token usage therefore come
short of a free reasoning agent on enumeration depth.
The gain comes at a price: a 5.5\,\% success rate (6 of 109 attempts) and roughly three times the
token cost (222,094 vs.\ 81,453).
The local \textit{qwen3:8b} is unusable on this scale, executing zero tool calls and producing no reconnaissance.


\section{Unexpected Results}\label{sec:discussion-unexpected}

First, the \textbf{unstructured-output configuration outperformed the structured
multi-agent setup on correctness} in the nslab (9 vs.\ 7).
The unstructured run surfaces all three severe findings \gls{SNMP}, open \gls{DNS}
recursion, and the cloned \gls{SSH} host keys.
This leads to the conclusion that the schema constraints which help with speed and reduce token usage come short in comparison
with a free reasoning agent.
The free reasoning agent issued UDP scans with additional TCP scans and executed NSE scripts that the multiagent specialised
scenario did not.
The results came with the price of a 5.5\,\% success rate and roughly three times the token cost (222{,}094 vs.\ 81{,}453).

Second, the \textbf{gain of the token reduction} from decomposition was larger
than expected.
A 68\,\% drop for \textit{qwen3:30b} shows that prompt-context growth, not raw
computation, dominates the cost of single-agent reconnaissance for smaller models.

Third, the \textbf{complete failure of the local \textit{qwen3:8b} model} (zero
\textit{cli\_tool} calls across all runs) was more absolute than anticipated.
The model did not perform any reconnaissance, falling back to a generic recommendation report or a report with connectivity issues.

Fourth, the \textbf{static scenario discovering more services yet fewer security
findings} than the \gls{AI} models.
This result was expected because of the complexity and reasoning that strengthen the claim to use ai models for security tasks.
The static scenario that performed rapidly in the virtual environment needed a lot of time to enumerate the network.


\section{Limitations}\label{sec:discussion-limitations}

\paragraph{Result asymmetry.}
As noted in Chapter~\ref{ch:evaluation}, the engineered scenario returns a structured
\textit{NetworkDiscoveryResultMap} while the \gls{AI} scenario, particularly in
unstructured mode, returns only a log/text stream.
A fully quantitative comparison therefore required either parsing the
free-form output or accepting this asymmetry.
This limits the precision of the findings of unstructured output.

\paragraph{Reliability and success rate.}
The headline figures report \(n = 10\) \emph{successful} runs, but the success factors
reveal unreliability.
In the cyberlab the success factor concludes 5 of 27 attempts for the multi agent and 6 of 109 runs for the unstructured attempts.
The runs were interrupted because of high token usage and missing comparison with other models, only the frontier model
could perform the tasks, though prone to errors.
The \gls{AI} scenario is not yet dependable enough for unsupervised operation, and the reported averages reflect a filtered best-case subset.

\paragraph{Incomplete environment and scenario matrix.}
The methodology (Section~\ref{subsec:environments}) foresaw three environments and a
blue-team intrusion-detection scenario (Section~\ref{subsubsec:intrusion-detection-blue-team}).
The evaluation covers the virtual Containerlab and the Cyberlab BFH segment but not a physical environment,
and the evaluation focuses on the red-team reconnaissance
scenario; the blue-team comparison (a \emph{should} goal) remains open.
In the process to implement \gls{AI} agents with LangChain, we spent a lot of time to enhance user experience and
implement a way to use structured output to evaluate the artifacts.
Therefore the should goals could not be reached and the physical environment uses the same devices and services as the
containerlab.
Therefore we invested more time in evaluating the other environments.

\paragraph{Subjectivity of qualitative scoring.}
The correctness and hallucination scores were assigned manually against the scale in
Table~\ref{tab:qualitative-scale}.
Although the scale is explicit, the scoring involves human judgement and is not
independently inter-rated, which introduces a degree of subjectivity into the
qualitative results.


\section{Recommendations}\label{sec:discussion-recommendations}

\begin{enumerate}
    \item \textbf{Match the model tier to the task.} Autonomous reconnaissance
    requires a capable model; the frontier \textit{claude-opus-4-7} is the only
    configuration that reliably extracted the Critical findings.
    Small local models such as \textit{qwen3:8b} are unsuitable for tool-driven
    reconnaissance and should not be relied upon.

    \item \textbf{Use multi-agent decomposition for smaller or self-hosted models.}
    Where a frontier model is unavailable for cost or data-locality reasons,
    task decomposition recovers much of the efficiency and consistency gap,
    most visibly the 68\,\% token reduction for \textit{qwen3:30b}.

    \item \textbf{Combine engineered breadth with \gls{AI} depth.} A hybrid pipeline:
    the deterministic scenario for exhaustive, reproducible service
    enumeration, followed by an \gls{AI} layer for semantic assessment and
    attack-chain reasoning, would pair the static scenario's coverage with
    the interpretive strength of the \gls{AI}.

    \item \textbf{Add a human-in-the-loop and retry strategy.} Given the low success
    rates, an operator checkpoint or automatic retry on incomplete runs would
    raise dependability to a level acceptable for practical use, and would
    mitigate the silent ``ping-sweep-only'' failures.
    While partly implemented in \todo{add section in AI implementation with middleware and dynamic tool call}
    we omitted the branch, because the agent always tried to use just the recommended tools and did not deliver a
    more sophisticated result.
\end{enumerate}
